\cotefile{12}
\scanpage{82}
$C$ courbe lisse \inserted{propre} sur $\mathbb{F}_{\unsure{q|p}}$

$A$ schéma abélien sur $C$.

Des conjectures de Weil, résulte que
\[
H^i(C,A)=0\qquad\text{si}\quad i\geq 2
\]
et sans doute, utilisant Tate, on doit avoir
\[
H^1(C,A)=0
\]

Soit $x$ un pt fermé de $C$, $U=C-\{x\}$, alors on a
\diagram{\scriptsize
\begin{tikzcd}[ampersand replacement=\&,column sep=.65em,row sep=.35em]
\deleted{\unsure{H^0(U,A)|H^0(C,A)}}\arrow[r] \&
\deleted{\unsure{H^1_x(C,A)}}\arrow[r] \&
H^1(C,A)\arrow[r] \& H^1(U,A)\arrow[r] \& H^2_x(C,A)\arrow[r] \& H^2(C,A)\\
\& \deleted{\unsure{0}} \& \shortparallel \& \& \& \shortparallel\\
\& \& 0 \& \& \& 0
\end{tikzcd}}
d'où
\[
H^1(U,A)\simeq H^2_x(C,A)
\]
Or par Tate,
\[
H^2_x(C,A)\deleted{\unsure{\simeq|=}}\;H^1(K_x,A)\deleted{\unsure{\simeq|=}}\quad
\text{\uline{dual} de }H^0(K_x,A^{*})\deleted{\unsure{\simeq}}
\]
\deleted{\unsure{ce|on} \unsure{n'est en général pas un} \unsure{groupe|groupe fini}\inserted{\unsure{groupe de}}} \uuline{c'est donc essentiellement un $p$-\deleted{\unsure{fini}}groupe}, modulo groupe fini ! O.K.

Peut se généraliser au cas où on prend $X$ régulier dont on enlève un diviseur régulier \textendash\,\textendash\,\textendash\,\textendash
% archive sheet number: 81
